% COURSE-ASSISTANT SOURCE — FALL 2025 ARCHIVE
% Historical administrative frames about contact details, Athena, lectures,
% assessment, seminars, and AI policy have been omitted. They must not be used
% as guidance for Fall 2026. Current-course material always takes priority.
% See LN00_figure_notes.md for descriptions of figures not stored here.
\documentclass[10pt,handout]{beamer}

\input{EC2211_preamble.txt}
\title{Course Introduction}

\begin{document}
\maketitle

\begin{frame}{Literature}
\begin{columns}
    \column{0.77\textwidth}
    \begin{itemize}
        \item Textbook: Charles I. Jones (2024), \textit{Macroeconomics}, W.W. Norton, 6E International Student Edition. \medskip
        \item Lecture notes (=presentation slides) and problem sets: These are important parts of the course. They also indicate what material and models you should focus on.\medskip
        \item Articles and policy reports: These are for policy perspectives or in-depth understanding. You should look more closely at articles marked by a *. Don't get stuck on technical material that we have not covered in class.
    \end{itemize}

    \column{0.23\textwidth}
    \begin{figure}
        \centering
        \includegraphics[width=0.8\linewidth]{JonesBook.jpg}
    \end{figure}
\end{columns}
\end{frame}



\begin{frame}{Charles "Chad" I. Jones}
\begin{columns}
    \column{0.6\textwidth}
    \begin{itemize}
        \item Professor, Stanford University (Graduate School of Business)\medskip
        \item Expert on economic growth, recently with a focus on artificial intelligence \medskip
        \item \href{https://web.stanford.edu/~chadj/}{https://web.stanford.edu/$\sim$chadj/}
    \end{itemize}

    \column{0.4\textwidth}
    \begin{figure}
        \centering
        \includegraphics[width=0.95\linewidth]{ChadinLafayette_2023tiny.jpg}
    \end{figure}
    \centering \small{Chad Jones}
\end{columns}
\end{frame}


\begin{frame}{Prerequisites}
\begin{itemize}
    \item I assume that you have taken an introductory course in macroeconomics, like EC1211 \textit{Macroeconomic Theory and Applications}\medskip
    \item If you have not, you will have to read the textbook more carefully as I will skip some material that I know was covered well in that course\medskip
    \item There will be some repetition of material covered in EC1211, but my ambition is that we can dig further into some interesting/important macroeconomic questions\medskip
\end{itemize}
\end{frame}


\begin{frame}{Contents}
\begin{itemize}
    \item Facts about economic aggregates and macroeconomic models\medskip
    \item Familiarize you with common macroeconomic tools and concepts, and develop your analytical skills\medskip
    \item Continuous references to historical and current macroeconomic phenomena\medskip
    \item Global perspective as well as special focus on Sweden\medskip
\end{itemize}
\end{frame}


\begin{frame}{What do macroeconomists do?}
\begin{itemize}
    \item[1.] Describe and measure the economy\smallskip
    \begin{itemize}
        \item National accounts, growth, unemployment, inflation, productivity, resource utilization, ...\smallskip
        \item Welfare, inequality\smallskip
        \item ...\smallskip
    \end{itemize}
    \item[2.] Analyze causality, policy options\smallskip
    \begin{itemize}
        \item Policies and institutions that promote economic growth\smallskip
        \item Economic effects of fiscal and monetary policy\smallskip
        \item ...\smallskip
    \end{itemize}
    \item[3.] Forecasting\smallskip
    \end{itemize}
\end{frame}


\begin{frame}{This course}
\begin{itemize}
    \item[1.] Focus on understanding how the economy works\smallskip
    \begin{itemize}
        \item Understand and interpret economic data\smallskip
        \item Understand economic relationships, causality\smallskip
    \end{itemize}
    \item[2.] We need models\smallskip
    \begin{itemize}
        \item Macroeconomists' laboratories\smallskip
    \end{itemize}
    \item[3.] Very little focus on forecasting\smallskip
    \item[4.] Much of our knowledge comes from economic failures and crises\smallskip
    \begin{itemize}
        \item We will consider specific current and historical events\smallskip
    \end{itemize}
\end{itemize}
\end{frame}


%%%%%%%%%%%%%%%%%%%%%%%%%%%%%%%%%%%%%%%%%%%%%%%%%%%
%%%%%%%%%%%%%%%%%%%%%%%%%%%%%%%%%%%%%%%%%%%%%%%%%%%
\section{Some examples}


\begin{frame}{What is GDP and what does it measure?}
\begin{figure}
    \centering
    \includegraphics[width=0.94\linewidth]{GDP_se.png}
    \flushright{\tiny{Current prices. Source: Statistics Sweden.}}
\end{figure}
\end{frame}


\begin{frame}{Why do some countries grow rapidly, others not?}
\begin{figure}
    \centering
    \includegraphics[width=0.96\linewidth]{Growth.png}
    \flushright{\tiny{Real GDP per capita, 2011 USD. Source: \href{https://www.rug.nl/ggdc/historicaldevelopment/maddison/releases/maddison-project-database-2023}{The Maddison Project}.}}
\end{figure}
\end{frame}


\begin{frame}{Why did inflation increase so much?}
\begin{figure}
    \centering
    \includegraphics[width=0.9\linewidth]{LN00b_Intro/HICP.png} \newline
    \flushright{\tiny{Percent annual change in HICP for Sweden and the Euro area, CPI for the United States. Sources: \href{https://ec.europa.eu/eurostat/databrowser/view/prc_hicp_midx__custom_18023313/default/table}{Eurostat} and \href{https://fred.stlouisfed.org/series/CPIAUCSL}{FRED}.}}
\end{figure}
\end{frame}


\begin{frame}{What can central banks achieve?}
\begin{figure}
    \centering
    \includegraphics[width=0.65\linewidth]{LN00b_Intro/FedDecision.png}
\end{figure}
\end{frame}


\begin{frame}{Is it important that central banks are independent?}
\centering
\begin{tikzpicture}
  % base image
  \node[inner sep=0] (base) {\includegraphics[width=0.74\linewidth]{Lect00a_Literature/TrumpFiresCook.png}};

  % top image, rotated and placed relative to base
  \only<2->{
  \node[inner sep=0, anchor=south east] 
        at ($(base.south east)+(-1.8cm,0.3cm)$)
        {\includegraphics[width=.3\linewidth, angle=-18, origin=c]{Figs/DJT.jpg}}
    };
\end{tikzpicture}\end{frame}


\begin{frame}{Should we worry about high and growing public debt?}
\begin{figure}
    \centering
    \includegraphics[width=0.9\linewidth]{LN00b_Intro/GovernmentDebt.png} \newline
    \flushright{\tiny{Source: \href{https://www.imf.org/en/Publications/WEO/weo-database/2025/april/select-country-group}{IMF World Economic Outlook}}}
\end{figure}
\end{frame}


%%%%%%%%%%%%%%%%%%%%%%%%%%%%%%%%%%%%%%%%%%%%%%%%%%%
%%%%%%%%%%%%%%%%%%%%%%%%%%%%%%%%%%%%%%%%%%%%%%%%%%%
\section{Models and data}


\begin{frame}{}
\begin{quote}
    {\large ``All theory depends on assumptions which are not quite true. That is what makes it theory. The art of successful theorizing is to make the inevitable simplifying assumptions in such a way that the final results are not very sensitive.''}
    \vskip5mm
    \hspace*\fill{\small--- Robert Solow}
\end{quote}
\end{frame}


\begin{frame}{Why models?}
\begin{itemize}
    \item Simplified versions of reality\medskip
    \item Models are economists' laboratories; once equilibrium is characterized we can conduct policy experiments\medskip
    \item Different models serve different purposes\medskip
    \item Impossible to incorporate all features in one model: tractability
\end{itemize}
\end{frame}


\begin{frame}{Modeling principles}
\begin{itemize}
    \item \alert{Exogenous variables}: treated as given in the analysis\smallskip
    \item \alert{Endogenous variables}: determined within the model\smallskip
    \item Variable of interest should be endogenously explained\smallskip
    \item The models analyze how exogenous variables affect the endogenous variables 
\end{itemize} \medskip \pause
\begin{figure}
    \centering
    \includegraphics[width=0.55\linewidth]{LN00b_Intro/MACRO6_FIG01.06.jpg}
\end{figure}
\end{frame}


\begin{frame}{Models: Agents}
\begin{itemize}
    \item Ideally, models should be based on the behavior of individual agents:
    \begin{itemize}
        \item Households maximize utility
        \item Firms maximize profits
    \end{itemize}\medskip
    \item Other agents and institutions:
    \begin{itemize}
        \item Financial institutions\smallskip
        \item Trade unions\smallskip
        \item Vested interest groups (political lobbies)\smallskip
        \item ...
    \end{itemize}
\end{itemize}
\end{frame}


\begin{frame}{Economic policy}
\begin{itemize}
    \item Economic policy: government may try to influence the general equilibrium by affecting the incentives of individual agents\medskip
    \item Government: fiscal authority + central bank\medskip
    \item Fiscal policy: taxes, subsidies, transfers, government consumption\medskip
    \item Monetary policy: interest rates, money supply, exchange rates
\end{itemize}
\end{frame}


\begin{frame}{Data}
\begin{itemize}
    \small{
    \item Data and econometric techniques used for:
    \begin{itemize}
        \item Descriptive analysis\smallskip
        \item Causal analysis (requires exogenous variation)\medskip
    \end{itemize}

    \item I will try to provide links to useful data sources along the course \medskip
    
    \item Some very useful sources for macro data:
    }
    \begin{itemize}
    \footnotesize{
        \item International organizations: \href{https://data-explorer.oecd.org/}{OECD Data Explorer}; \href{https://www.imf.org/en/Publications/SPROLLs/world-economic-outlook-databases}{IMF World Economic Outlook}; \href{https://data.worldbank.org/}{the World Bank}; \href{https://www.bis.org/statistics/dataportal/}{BIS}\smallskip
        \item Statistical agencies: \href{https://www.statistikdatabasen.scb.se/pxweb/en/ssd/}{Statistics Sweden}; \href{https://ec.europa.eu/eurostat/data/database}{Eurostat}; \href{https://www.bls.gov/data/}{Bureau of Labor Statistics} \smallskip
        \item Private initiatives: \href{https://www.rug.nl/ggdc/productivity/pwt/}{Penn World Tables} (annual macro data 1950-2019 for 183 countries; \href{https://www.rug.nl/ggdc/historicaldevelopment/maddison/releases/maddison-project-database-2023}{the Maddison Project} (very long time series on growth and income)  \smallskip
        \item Central banks: \href{https://www.riksbank.se/en-gb/statistics/}{Riksbank} (e.g. interest rates and exchange rates); \href{https://data.ecb.europa.eu/data/datasets}{ECB} (a lot of European financial data); \href{https://fred.stlouisfed.org}{St Louis Fed's FRED database} (mostly U.S. macro and financial data, very user-friendly!)
        }
    \end{itemize}
\end{itemize}
\end{frame}



\end{document}


